\section{Related Work}
\label{sec:related}

\subsection{Android malware detection and classification}

Early Android malware detection was dominated by signature matching and
manifest-level permission analysis~\cite{felt2011survey}. As
obfuscation techniques matured, work shifted to
decompiled-code features~\cite{arp2014drebin} and, more recently, to
dynamic behavioural traces captured in instrumented
sandboxes~\cite{ali2020ransomware}. The availability of corpora such as
Drebin, CICAndMal2017, and UMD has accelerated comparative benchmarking
of classical and deep models.

\subsection{Markov chains and associative rules for malware}

Markov models have a long history in malware classification: sequences of
system calls, API calls, or opcodes are modelled as first-order or
$k$-order transitions, and the resulting transition probabilities serve
as features for downstream classifiers. D'Angelo et
al.~\cite{dangelo2023association} recently extended this paradigm by
mining $k$-spaced \emph{associative rules} between API calls, pruning
by support and confidence, and classifying via per-class softmax ranking
(Eq.~7 of their paper). Their federated-learning variant yields privacy
guarantees at a small accuracy cost; we use their centralised
formulation as our primary statistical baseline.

\subsection{Sequence models and transformers}

Transformers~\cite{vaswani2017attention} revolutionised natural language
processing and have since been applied to system-call sequences for
malware detection with strong results. However, standard softmax
attention creates a linear composition between the value projection
$V$ and the output projection $W_O$ — a low-rank bottleneck — and
empirically produces ``attention sinks'' where early tokens receive
high weight regardless of input~\cite{qiu2025gated}.

\subsection{Gated attention}

Qiu et al.~\cite{qiu2025gated} introduced a family of gating schemes
for large language model attention. Among three positions they
evaluate (G1 after SDPA output, G2 on V before SDPA, G3 after $W_O$),
the G1 position — a head-specific sigmoid gate between the scaled
dot-product attention output and the output projection — was
consistently best. It introduces non-linearity, input-dependent
sparsity, and removes the attention-sink artefact.
We adopt the G1 formulation and initialise the gate bias to
$-2.0$ following Qiu et al.; we observe a different converged
sparsity regime on our task (Section~\ref{sec:results}).

\subsection{Explainable malware classification}

Post-hoc attribution methods such as SHAP~\cite{lundberg2017shap} and
LIME have been applied to malware classifiers to rank features by
importance. These methods fit a surrogate model locally around each
prediction; the resulting explanations are approximations, not
exactly the reasoning of the underlying classifier. In contrast, our
gate activations are produced by the same forward pass as the
classification decision, so they are \emph{intrinsic} to the model's
computation and do not require a surrogate.
